\section*{Methods}

\subsection*{Study Area}
Foxes were captured and collared in four different areas along a gradient of
decreasing landscape productivity and human land use as latitude increased from
58$^{\circ}$ to 62$^{\circ}$ N in central Sweden and Norway. The southernmost
area consists of fragmented boreonemoral forests, agricultural lands, and human
settlements, whereas the northern areas are dominated by boreal forests and
scattered human settlements. Norway spruce \emph{Picea abies}, Scots pine
\emph{Pinus sylvestris}, and to a lesser degree birch \emph{Betula} spp.\ are
the most common tree species. Broad-leaved deciduous species such as English oak
\emph{Quercus robur} and Norway maple \emph{Acer platanoides} are present in
the southern area but rare in the north. The northernmost part grades into
alpine tundra of low diversity and productivity. Altitude increases with
latitude from approximately 25 to 750 m above sea level, and snow cover shifts
from irregular in the south to continuous from November through May in the
north. For further details see \citet{waltonVariationHomeRange2017}.

\subsection*{Capture, Tagging \& Tracking}
All captured foxes heavier than 5 kg were fitted with a GPS radio collar
(Tellus Ultralight, 210 g, Televilt Inc., Lindesberg, Sweden). Foxes were
sexed, measured, weighed, and aged at capture. Age was classified as sub-adult
($< 1$ year) or adult ($> 1$ year) based on the degree of tooth wear and tooth
coloration. GPS collars were programmed to acquire a minimum of three positions
per day; variation in duty cycles resulted in collar drop-off after 6--9
months. All capture and handling procedures were approved by the Swedish Animal
Ethics Committee (permit numbers DNR 70--12, DNR 58--15) and the Norwegian
Experimental Animal Ethics Committee (permit numbers 2009/122825, 2012/20038,
2014/207803). Permits to capture wild animals were issued by the Norwegian
Directorate for Nature Management and the Swedish Environmental Protection Board
(NV-03459-11). For further details see
\citet{waltonVariationHomeRange2017,waltonLongdistanceDispersalRed2018}.

Between 2011 and 2019, we captured and collared 132 individual red foxes.
One fox (Ann) was tracked continuously using two successive collars and its
records were merged into a single uninterrupted interval. Five foxes (Frans,
Gunilla, Ingvar, Oskar, Cristina) were recaptured after a gap in monitoring and
treated as two independent observation intervals, consistent with the
Andersen-Gill counting-process framework \citep{andersenCoxsRegressionModel1982}.
Cristina was sub-adult at first capture and adult at recapture, and her two
intervals were assigned the appropriate age class separately. We excluded four
individuals (Gunilla's first capture, Källsvedan, Oskar's second capture, and
Stiga) that contributed fewer than two weeks of tracking and were therefore
uninformative for survival estimation, leaving a total of 127 individual
fox-intervals in the analysis (Table~\ref{table:capt}).

An inactive GPS collar transmitted a mortality signal and SMS notification used
to record the approximate time of death. Mortality causes other than hunting
were determined in the field: damage consistent with vehicle strikes was
assessed visually; carcasses exhibiting alopecic patches, crusted skin, and
sometimes a foul aromatic odour were classified as sarcoptic mange; and a wolf
kill was identified by tracks and bite marks at the kill site. Most mortalities
resulted from hunting, and hunters returned collars with the date of harvest.
Cases where cause of death could not be determined were classified as unknown.

\begin{table}[t]
\centering
\captionsetup{font=small}
\caption{\textbf{Number of collared fox intervals included in the analysis,
  by age class and sex. One fox (Ann) contributed as a single merged interval
  across two collars; five foxes contributed two intervals each following
  recapture. Four individuals tracked for fewer than two weeks were
  excluded.}}
\label{table:capt}
\begin{tabular}{lrr|r}
\toprule
Sex     & Adult & Sub-adult & Total \\
\midrule
Females & 28    & 22        & 50    \\
Males   & 44    & 33        & 77    \\
\midrule
Total   & 72    & 55        & 127   \\
\bottomrule
\end{tabular}
\end{table}

\subsection*{Analysis}
To estimate the effects of age, sex, and season on survival we used the
Andersen-Gill extension of the Cox proportional hazards (CPH) regression model
\citep{andersenCoxsRegressionModel1982}, following approaches used for similar
telemetry studies \citep{johnsonModelingSurvivalApplication2004,
ahlenSurvivalFemaleCapercaillie2013}. The non-parametric Kaplan-Meier estimator
was used to visualise and inspect survival curves
\citep{pollockSurvivalAnalysisTelemetry1989}.

Each fox contributed one row per calendar week it was present in the study,
yielding a counting-process dataset in which each row represents the interval
$(t_{\text{start}}, t_{\text{end}}]$ for that individual-week. The last week a
fox was monitored was recorded as either a known mortality or a censored
observation due to collar drop-off or failure; all preceding weeks were
censored. Collar failure and collar drop-off (CAUSE = 9) were treated as
non-informative right-censoring throughout. We pooled study areas and years due
to low overall sample size.

We set the start of the red fox year to ISO week 36 (early September), when
cubs typically separate from family groups, and expressed time within the year
as a continuous fox-week variable running from 1 (week 36) to 52 (week 35 of
the following calendar year). Sub-adults whose tracking extended beyond week 35
of the calendar year following their capture year were reclassified as adults
from that point onward; two individuals (Cristina and Gijom) met this criterion.

Robust sandwich standard errors, obtained via \texttt{cluster(FoxID)} in the
\texttt{coxph} function, accounted for the non-independence of multiple
observation rows from the same individual
\citep{therneauPackageSurvivalAnalysis2024}. We developed and compared five
candidate CPH models containing combinations of age (sub-adult vs.\ adult),
sex, and season (Table~\ref{table:models}). Season was evaluated as: (i) a
binary variable (autumn/winter vs.\ spring/summer), (ii) a continuous linear
fox-week covariate, (iii) a sine-cosine pair capturing circular annual
periodicity, and (iv) a four-level factor (autumn, winter, spring, summer).
Models were compared using Akaike's Information Criterion (AIC) computed from
the partial log-likelihood. Likelihood ratio tests (LRT) between nested models
were performed on refitted models without the \texttt{cluster()} argument,
because the robust variance estimator does not alter the partial log-likelihood
and therefore does not affect LRT statistics. We assessed whether a sex
$\times$ season interaction improved model fit, but given only nine female
mortality events and the rule of thumb of approximately ten events per
estimated parameter, this interaction was considered underpowered and is
reported for transparency only.

The proportional hazards assumption was evaluated using scaled Schoenfeld
residuals via \texttt{cox.zph} with a rank time transformation
\citep{therneauPackageSurvivalAnalysis2024}. All data preparation and analyses
were conducted in R version 4.x \citep{rcoreteamLanguageEnvironmentStatistical2024},
using the packages \texttt{survival}
\citep{therneauPackageSurvivalAnalysis2024} and \texttt{survminer}
\citep{kassambaraSurvminerDrawingSurvival2024}.

\begin{table}[t]
\centering
\captionsetup{font=small}
\caption{\textbf{Number of weeks contributed to the analysis by sex and age
  class across all study areas and years (127 fox-intervals, 35 known
  mortalities).}}
\label{table:weeks}
\begin{tabular}{llrr}
\toprule
Sex     & Age class  & Weeks & Mortalities \\
\midrule
Males   & Adult      & 936   & 11 \\
Males   & Sub-adult  & 663   & 15 \\
Females & Adult      & 650   &  4 \\
Females & Sub-adult  & 411   &  5 \\
\midrule
Total   &            & 2660  & 35 \\
\bottomrule
\end{tabular}
\end{table}
